\documentclass[11pt]{article}
\usepackage[margin=1in]{geometry}
\usepackage{amsmath,amssymb,amsthm}
\usepackage[T1]{fontenc}
\usepackage{lmodern}
\usepackage{hyperref}

\newtheorem{theorem}{Theorem}
\newtheorem{lemma}{Lemma}
\newtheorem{corollary}{Corollary}

\newcommand{\Jal}{J_{\aleph}}
\newcommand{\Ree}{\operatorname{Re}}
\newcommand{\dens}{\operatorname{dens}}

\title{A cardinal Jung dichotomy for spaces \(C(K)\)}
\author{Agent lane 11}
\date{2026-07-03}

\begin{document}
\maketitle

\section*{Source and Target}

The source paper is J. M. F. Castillo and P. L. Papini,
\emph{The separable Jung constant in Banach spaces}, arXiv:1910.02402.
In the subsection ``On the stability of the Jung constants'', the authors
ask whether, for a space \(C(K)\) of continuous functions on a compact
space, the inequality \(J_{\aleph}(C(K))<2\) implies
\(J_{\aleph}(C(K))=1\).

This packet proves that question for every uncountable cardinal \(\aleph\).
The exact result is the dichotomy
\[
        J_{\aleph}(C(K))\in\{1,2\},
\]
with the value determined by whether \(K\) is an \(F_{\aleph}\)-space.

\section*{Definitions}

For a bounded set \(A\) in a Banach space \(X\), write
\[
 \delta(A)=\sup\{\|a-b\|:a,b\in A\},\qquad
 r(A)=\inf_{x\in X}\sup_{a\in A}\|a-x\|.
\]
For an uncountable cardinal \(\aleph\), the cardinal Jung constant is
\[
 J_{\aleph}(X)=\sup \frac{2r(A)}{\delta(A)},
\]
where the supremum is over closed bounded sets \(A\subset X\) with
\(\delta(A)>0\) and \(\dens(A)<\aleph\).

A compact Hausdorff space \(K\) is an \(F_{\aleph}\)-space if every two
disjoint open subsets of \(K\) which are unions of fewer than \(\aleph\)
closed sets have disjoint closures.

\section*{Proof Intuition}

The positive direction is already encoded in the cited injectivity theory:
\(J_{\aleph}(X)=1\) means \((1,\aleph)\)-injectivity, and for \(C(K)\)
spaces that is exactly the \(F_{\aleph}\) property of \(K\).

For the converse, the failure of \(F_{\aleph}\) supplies two disjoint open
sets \(G,H\), built from fewer than \(\aleph\) closed pieces, whose closures
touch. Continuous functions supported on the closed pieces of \(G\) demand
that any small-radius center be positive on \(\overline G\); the negative
of the analogous functions on \(H\) demands that it be negative on
\(\overline H\). At a common closure point these demands are incompatible.

\begin{lemma}\label{lem:non-faleph-witness}
Let \(\aleph\) be an uncountable cardinal. If \(K\) is not an
\(F_{\aleph}\)-space, then \(C(K)\) contains a closed bounded set \(A\) with
\(\dens(A)<\aleph\), \(\delta(A)=1\), and \(r(A)=1\).
\end{lemma}

\begin{proof}
Since \(K\) is not an \(F_{\aleph}\)-space, there are disjoint open sets
\[
        G=\bigcup_{i\in I} C_i,\qquad
        H=\bigcup_{j\in J} D_j
\]
with \(|I|,|J|<\aleph\), where each \(C_i,D_j\) is closed in \(K\), and
\[
        \overline G\cap\overline H\neq\varnothing .
\]
For each \(i\in I\), Urysohn's lemma gives \(u_i\in C(K)\) with
\[
        0\leq u_i\leq 1,\qquad u_i|_{C_i}=1,\qquad u_i|_{K\setminus G}=0.
\]
For each \(j\in J\), choose \(v_j\in C(K)\) with
\[
        0\leq v_j\leq 1,\qquad v_j|_{D_j}=1,\qquad v_j|_{K\setminus H}=0.
\]
Set
\[
        A_0=\{u_i:i\in I\}\cup\{-v_j:j\in J\}\subset C(K),
\]
and let \(A=\overline{A_0}\) in \(C(K)\). Since \(|I|+|J|<\aleph\), the
set \(A_0\) has cardinal \(<\aleph\), and therefore \(\dens(A)<\aleph\).

All \(u_i\)'s and \(v_j\)'s take values in \([0,1]\), and the supports of
the \(u_i\)'s are disjoint from the supports of the \(v_j\)'s. Hence
\[
 \|u_i-u_{i'}\|\leq 1,\qquad
 \|-v_j+v_{j'}\|\leq 1,\qquad
 \|u_i-(-v_j)\|=\|u_i+v_j\|\leq 1.
\]
Thus \(\delta(A)\leq 1\). Since \(G\) and \(H\) are nonempty, choose
\(i\in I\) with \(C_i\neq\varnothing\) and \(j\in J\). At any
\(x\in C_i\) we have \(u_i(x)=1\) and \(v_j(x)=0\), so
\(\|u_i-(-v_j)\|\geq 1\). Therefore \(\delta(A)=1\).

The zero function is a center of radius \(1\), so \(r(A)\leq 1\). Suppose,
toward a contradiction, that \(h\in C(K)\) and \(\rho<1\) satisfy
\[
        \sup_{a\in A}\|h-a\|\leq \rho .
\]
For every \(i\in I\) and every \(x\in C_i\),
\[
        \Ree h(x)\geq 1-\rho,
\]
because \(u_i(x)=1\). Hence \(\Ree h\geq 1-\rho\) on \(G\), and by
continuity also on \(\overline G\). Similarly, for every \(j\in J\) and
every \(y\in D_j\),
\[
        \Ree h(y)\leq -1+\rho,
\]
because \(-v_j(y)=-1\). Hence \(\Ree h\leq -1+\rho\) on
\(\overline H\). At any point
\(p\in\overline G\cap\overline H\), this gives
\[
        1-\rho\leq \Ree h(p)\leq -1+\rho,
\]
which is impossible for \(\rho<1\). Thus \(r(A)\geq 1\), and so
\(r(A)=1\).
\end{proof}

\begin{theorem}\label{thm:dichotomy}
Let \(\aleph\) be an uncountable cardinal. For every compact Hausdorff
space \(K\),
\[
        J_{\aleph}(C(K))=
        \begin{cases}
        1, & K\text{ is an }F_{\aleph}\text{-space},\\
        2, & K\text{ is not an }F_{\aleph}\text{-space}.
        \end{cases}
\]
\end{theorem}

\begin{proof}
If \(K\) is an \(F_{\aleph}\)-space, then the characterization of
Aviles--Cabello Sanchez--Castillo--Gonzalez--Moreno says that \(C(K)\) is
\((1,\aleph)\)-injective. Castillo--Papini prove that a Banach space \(X\)
is \((1,\aleph)\)-injective if and only if \(J_{\aleph}(X)=1\). Hence
\(J_{\aleph}(C(K))=1\).

If \(K\) is not an \(F_{\aleph}\)-space,
Lemma~\ref{lem:non-faleph-witness} gives a closed bounded set
\(A\subset C(K)\) with \(\dens(A)<\aleph\), \(\delta(A)=1\), and
\(r(A)=1\). Therefore
\[
        J_{\aleph}(C(K))\geq \frac{2r(A)}{\delta(A)}=2.
\]
For every Banach space \(X\), \(J_{\aleph}(X)\leq 2\), since every
bounded set has radius at most its diameter by centering at one of its own
points. Consequently \(J_{\aleph}(C(K))=2\).
\end{proof}

\begin{corollary}
For every uncountable cardinal \(\aleph\) and compact Hausdorff space
\(K\), if \(J_{\aleph}(C(K))<2\), then \(J_{\aleph}(C(K))=1\).
\end{corollary}

\section*{Verification and Novelty Check}

The cheap run indexes were searched for the arXiv id, the paper title, and
the phrases \(J_s(C(K))\), \(J_{\aleph}(C(K))\), separable Jung constant,
Jung constant \(C(K)\), \(F\)-space, and \(F_{\aleph}\)-space. The only
direct hit for this target was this lane's earlier separable packet, now
upgraded in place.

The local source of arXiv:1406.6733 was checked for the definition of
\(F_{\aleph}\)-spaces and the characterization of \((1,\aleph)\)-injective
\(C(K)\) spaces. Web searches on 2026-07-03 for exact phrases including
``\(J_{\aleph}(C(K))\)'', ``\(F_{\aleph}\) Jung constant \(C(K)\)'', and
``The separable Jung constant in Banach spaces \(J_{\aleph}\)'' found the
source paper and the cited cardinal-injectivity paper, but no close exact
resolution.

\section*{References}

\begin{itemize}
\item J. M. F. Castillo and P. L. Papini, \emph{The separable Jung constant
in Banach spaces}, arXiv:1910.02402.
\item A. Aviles, F. Cabello Sanchez, J. M. F. Castillo, M. Gonzalez, and
Y. Moreno, \emph{\(\aleph\)-injective Banach spaces and
\(\aleph\)-projective compacta}, Rev. Mat. Iberoam. 31 (2015), 575--600;
arXiv:1406.6733.
\end{itemize}

\end{document}
